%! Unit 1, Lessons 1.0–1.4 — Quiz Study Guide (reference sheet; no answer key by design)
% This is a REFERENCE sheet, not a worksheet: every definition, every worked move, and every
% trap is already filled in, so there is nothing to answer and therefore no _key sibling.
% The practice a student does before the quiz is ../sample_quiz/, which HAS a key.
% NAMESTRIP: no \namedateperiod — the name row lives on a packet cover and the unit tests only.
% Vocabulary below is quoted from the five lessons' own vocabulary boxes (notes_key /
% experience_key \termans entries) so a student reading this reads the same words twice.
\documentclass[12pt]{article}
\usepackage{apstats-article}
\usepackage{apstats-boxes}
\usepackage{ltablex}
\keepXColumns
% One tabularx per group of lessons rather than one long table: a tabularx with X
% columns is an unbreakable box, so a single 25-row table floats whole to a fresh page
% and strands most of page 1. Three short tables flow, and each carries its lesson tag.

\newcommand{\sghead}[1]{%
  \vspace{0.13in}
  \boxguard[10]
  \begin{headlinebox}{navy}\color{white}\bfseries\small #1\end{headlinebox}
  \vspace{0.09in}}

\begin{document}

\pageheader{Unit 1 \quad Lessons 1.0--1.4}{Quiz Study Guide}

\vspace{0.12in}

\boxguard[8]
\begin{remindbox}
\small
The quiz covers \textbf{Lesson 1.0} through \textbf{Lesson 1.4} --- the vocabulary of data,
categorical vs.\ quantitative variables, displays for each, and describing a distribution.
Nothing on it needs a calculation harder than a percent. What it does need, every single time,
is an answer written \textbf{in context}: name the variable and its units, and say \emph{how}
you know. An answer that would be true of any data set anywhere earns nothing.
\end{remindbox}

\sghead{1. Vocabulary you are expected to use, not just recognize}

\small
\renewcommand{\arraystretch}{1.26}

\textbf{From Lessons 1.0 and 1.1 --- what data is, and the two kinds of variable}
\par\vspace{3pt}
\noindent\begin{tabularx}{\linewidth}{@{} l X @{}}
\rowcolor{sky}
\textbf{Term} & \textbf{What it means} \\
\midrule
Individual & one person, animal, or object described by \emph{one row} of the table. \\
Variable & a characteristic recorded for every individual --- a \emph{column heading}. \\
Value & one entry under a variable for one individual. \\
Variability & values differ from one individual to the next. This is why statistics exists. \\
Statistical question & anticipates variability \emph{and} is answered by collecting data. \\
Categorical variable & its values are category names or group labels --- Produce, Yes, Cedar Falls. \\
Quantitative variable & its values record a measured or counted \emph{amount} --- minutes, pounds, dollars. \\
The amount test & ask whether the value records \emph{how much or how many} --- not whether it
                   is written in digits. \\
Grouping (binning) & replacing measured amounts with group labels. Turns quantitative into
                   categorical, and the individual amounts cannot be recovered. \\
\end{tabularx}

\vspace{7pt}
\textbf{From Lessons 1.2 and 1.3 --- displaying each kind}
\par\vspace{3pt}
\noindent\begin{tabularx}{\linewidth}{@{} l X @{}}
\rowcolor{sky}
\textbf{Term} & \textbf{What it means} \\
\midrule
Frequency table & the \emph{count} in each category; the counts add to the group size. \\
Relative frequency table & the \emph{share} in each category; the entries add to 1, or 100\%. \\
Bar graph & one bar per category; gaps between the bars. \\
The comparison rule & to compare groups of \textbf{different sizes}, use relative frequencies ---
                   never raw counts. \\
Discrete variable & a count; nothing sits between two neighbouring values. \\
Continuous variable & a measurement; between any two values there is always another. \\
Histogram & grouped quantitative data, one bar per interval, height = count,
                   \textbf{bars touching}. \\
Interval width & the size of each histogram interval. Changing it changes the picture without
                   changing a single data value. \\
\end{tabularx}

\vspace{7pt}
\textbf{From Lesson 1.4 --- describing what the display shows}
\par\vspace{3pt}
\noindent\begin{tabularx}{\linewidth}{@{} l X @{}}
\rowcolor{sky}
\textbf{Term} & \textbf{What it means} \\
\midrule
Symmetric & the left half is roughly a mirror image of the right half. \\
Skewed right / left & named for the side the \textbf{longer tail} runs out on --- right: a few
                   values much larger; left: a few much smaller. \\
Unimodal / bimodal / uniform & one main peak; two clear peaks; all bars about equal, no peak. \\
Outlier & a value unusually far from the rest --- \textbf{investigated}, never simply deleted. \\
Gap / cluster & a stretch where nothing was observed; a bunch of values set off by a gap. \\
\end{tabularx}

\normalsize

\sghead{2. The five moves the quiz will ask you to make}

\small
\textbf{Move 1 --- Read a data table (1.0).} Rows are the individuals, column headings are the
variables, entries are the values. \emph{Worked:} in a table with one row per market vendor and
columns \emph{Stall type} and \emph{Weekly revenue}, the individuals are the \textbf{vendors},
the variables are \textbf{stall type} (categorical) and \textbf{weekly revenue} (quantitative),
and ``Produce'' is a \textbf{value}, not a variable.

\vspace{5pt}
\textbf{Move 2 --- Classify a variable (1.1).} Run the amount test, then --- if quantitative ---
ask counted or measured. \emph{Worked:} \emph{Stall number 101--148} is written in digits but
records no amount, so it is \textbf{categorical}. \emph{Bottles sold} is a count, so
\textbf{discrete quantitative}. \emph{Pounds of tomatoes} is a measurement, so
\textbf{continuous quantitative} --- and it stays continuous even if every value was rounded to a
whole pound.

\vspace{5pt}
\textbf{Move 3 --- Build and use the two tables (1.2).} Counts add to the group size; relative
frequencies add to 100\%. \emph{Worked:} 18 of 48 vendors sell produce, so the relative frequency
is $18/48 = 0.375$, or 37.5\%. To compare that market with one where 15 of 30 vendors sell
produce ($15/30 = 50\%$), you must use those shares: 18 is the larger \emph{count} and 50\% is
the larger \emph{share}, and the question asks which.

\vspace{5pt}
\textbf{Move 4 --- Read a quantitative display (1.3).} A dotplot and a stem plot give back every
individual value; a histogram does not --- it gives back only how many landed in each interval.
Histogram bars \textbf{touch} because the intervals run into one another; bar-graph bars have
gaps because the categories do not.

\vspace{5pt}
\textbf{Move 5 --- Describe a distribution (1.4).} All four, every time, in context:
\textbf{shape}, \textbf{unusual features}, \textbf{center}, \textbf{variability}.
\normalsize

\vspace{0.10in}
\boxguard[14]
\begin{scenariobox}[The sentence that earns full credit]{navy}
\small
``Weekly revenue for the 30 vendors is \textbf{skewed right}, with a \textbf{gap} between
\$1{,}000 and \$1{,}900 and one \textbf{outlier} near \$1{,}900. The \textbf{center} is about
\$500 per week, and the values run from about \$200 to \$1{,}900, so there is a great deal of
\textbf{variability} among the vendors.''
\par\vspace{4pt}
Four things. The variable named. The units named. If you can delete the context from your
sentence and still have it make sense, it is not yet an answer.
\end{scenariobox}

\sghead{3. Five traps that cost points every year}

\small
\begin{enumerate}[leftmargin=1.6em, itemsep=5pt]
  \item \textbf{Digits are not the test.} ZIP codes, jersey numbers, and stall numbers are
        categorical. Ask what the value \emph{records}, not what it looks like.
  \item \textbf{Skew is named for the tail, not the pile.} Most of the dots on the left with a
        few stragglers reaching right is \textbf{skewed right}.
  \item \textbf{A center by itself describes nothing.} Two checkout lines can share a middle wait
        of 4 minutes and be nothing alike. Always pair center with variability.
  \item \textbf{Different-sized groups need shares.} A bigger count out of a bigger group can be
        the smaller share. Say which number you used and why.
  \item \textbf{An outlier is a finding, not an error.} Report it, investigate it, and describe
        the distribution with it. Never delete a value because it is inconvenient.
\end{enumerate}
\normalsize

\sghead{4. Before you walk in}

\small
Work the \textbf{sample quiz} without your notes first, then check it against its key --- that
order is the whole point of it. If a free-response answer of yours is shorter than the key's, it
is almost certainly missing the context, not the arithmetic. Re-read the \textbf{Watch For} idea
from each lesson's cover, and be ready to say, out loud, what a histogram throws away.
\normalsize

\end{document}
